\originalchapter{尺度假设、有限尺寸与重整化}
\label{ch:scaling}
\source{8.334 L6--L8。}
\originalsection{自由能的齐次结构}
以$t=(T-T_c)/T_c$为温度尺度场、$h$为外场，假设奇异自由能在稳定临界固定点附近满足$f_s(t,h)=b^{-d}f_s(tb^{y_t},hb^{y_h})$。这个齐次形式来自粗粒化后每个相关块贡献有限自由能及尺度场的线性流，是需要临界尺度区的假设，不是任意函数恒等式。

选$b=|t|^{-1/y_t}$，得$f_s=|t|^{d/y_t}\Phi_\pm(h/|t|^{y_h/y_t})$。定义$\xi\sim|t|^{-\nu}$、$M\sim(-t)^\beta$、$\chi\sim|t|^{-\gamma}$、$C_s\sim|t|^{-\alpha}$、临界$M\sim h^{1/\delta}$，对该自由能求导得到
\begin{equation}
\nu=\frac1{y_t},\quad\alpha=2-d\nu,\quad\beta=\nu(d-y_h),\quad\gamma=\nu(2y_h-d),\quad\delta=\frac{y_h}{d-y_h}.
\end{equation}
临界关联$G(x)\sim|x|^{-(d-2+\eta)}$给场尺度$x_\phi=(d-2+\eta)/2$，源项$\int h\phi$无量纲要求$y_h=d-x_\phi$，从而$\gamma=\nu(2-\eta)$。这些推导解释指数关系为何不是独立猜测。
\originalsection{假设边界与危险无关变量}
$\alpha=2-d\nu$称超尺度关系，要求只有$\xi$控制奇异自由能体积尺度。在$d>4$的$\phi^4$，$u$在RG中流向0，却有序相$\phi_0\sim\sqrt{|r|/u}$，不能在求振幅前置$u=0$。这种危险无关变量使通常超尺度失效，平均场$\nu=1/2,\alpha=0$不会在任意$d$满足$2-d\nu$。边际维数还可有对数；BKT的$\xi$为本质奇异，也不能硬套普通温度幂指数。
\originalsection{RG变换与相关、无关方向}
先积分掉尺度$a$到$ba$之间的短波，再把坐标缩回原截止，最后重标场。所有对称允许的耦合组成向量$g$，变换$g'=R_b(g)$。固定点$R_b(g_*)=g_*$，在其附近$\delta g_i'=\sum_jM_{ij}\delta g_j$，对角尺度方向的本征值$b^{y_i}$。$y_i>0$相关、$<0$无关、$=0$边际，边际需看非线性项。对相关方向的名字应与所选基相关，原始温度可混入其他参数。

多个微观模型流到同一固定点，就有相同长距离指数与经尺度单位归一的普适函数；“所有材料都一样”并非普适性的含义。对称、维数、相互作用范围、守恒与拓扑等改变可使流到不同固定点。

Gaussian模型在$x=bx'$、$\phi(x)=b^{-(d-2)/2}\phi'(x')$下保持梯度项，质量$r'=b^2r$，四次耦合$u'=b^{4-d}u$。这直接给$y_r=2,y_u=4-d$与上临界维数。后续相互作用RG将这些工程维数修正。
\originalsection{有限尺寸标度}
当样品边长$L$限制相关长度，临界量依赖$tL^{1/\nu}$。例如$M=L^{-\beta/\nu}\mathcal M(tL^{1/\nu},hL^{y_h})$、$\chi=L^{\gamma/\nu}\mathcal X(tL^{1/\nu})$。Binder比$U_L=1-\avg{M^4}/[3\avg{M^2}^2]$无量纲，在相同几何边界下临界曲线近交于普适值。有限尺寸修正$L^{-\omega}$、外形及边界都会移动交点，不能用两条小样品曲线交点声称精确$T_c$。

有限尺寸谱或模拟数据做塌缩时，应先固定可检验的指数、误差与修正方案，再检查不同尺寸是否一致。任意调横纵坐标直到看起来重合不构成独立指数验证。
